% ---
% title: mooc题目
% date: 2026-07-01
% author: Crystal-Sky
% tags: 创新实践, MOOC, 题目整理
% summary: 《创新实践与展示》MOOC 题目与答案整理，依据 PDF 重新整理，包含单选题、多选题、判断题和填空题。
% slug: innovation-practice-mooc-questions
% pdf: assets/mooc-questions.pdf
% ---

\begin{abstract}
这是一份《创新实践与展示》MOOC 题目与答案整理，依据 PDF 文档重新提取并排版，按单选题、多选题、判断题和填空题归类。
\end{abstract}

\section{一、单选题}

\begin{question}{1. 下面哪一项数据结构不允许有重复元素？}
\begin{enumerate}[label=\Alph*., leftmargin=2.4em, itemsep=0.2em]
\item tuple
\item list
\item dict
\item set
\end{enumerate}
\end{question}
\begin{answer}
D（set）
\end{answer}
\begin{explanation}
set 是集合类型，集合天然去重；list、tuple 可以保存重复元素，dict 的键不能重复但题目问的数据结构通常指集合不允许重复元素。
\end{explanation}

\begin{question}{2. x=4，x*=(x+3)//2，经过上边两条语句运算后，x 的值是多少？}
\begin{enumerate}[label=\Alph*., leftmargin=2.4em, itemsep=0.2em]
\item 12
\item 16
\item 10
\item 8
\end{enumerate}
\end{question}
\begin{answer}
A（12）
\end{answer}
\begin{explanation}
先算右边：(x+3)//2=(4+3)//2=7//2=3；x*=3 等价于 x=x*3，所以结果为 12。
\end{explanation}

\begin{question}{3. PWM 的占空比是（）}
\begin{enumerate}[label=\Alph*., leftmargin=2.4em, itemsep=0.2em]
\item 高电平时间与周期时间的比值
\item 低电平持续的时间
\item 高电平持续的时间
\item 低电平时间与周期时间的比值
\end{enumerate}
\end{question}
\begin{answer}
A（高电平时间与周期时间的比值）
\end{answer}
\begin{explanation}
占空比定义为一个周期内高电平持续时间与整个周期时间之比。
\end{explanation}

\begin{question}{4. 假如 PCF8591 的基准源是 5.1V，那么它的分辨率是多少？}
\begin{enumerate}[label=\Alph*., leftmargin=2.4em, itemsep=0.2em]
\item 0.04
\item 0.02
\item 0.03
\item 0.01
\end{enumerate}
\end{question}
\begin{answer}
B（0.02）
\end{answer}
\begin{explanation}
PCF8591 是 8 位 A/D，量化级数约为 256，分辨率约为 5.1/256≈0.0199V，约等于 0.02V。
\end{explanation}

\begin{question}{5. 下列数中最大的是（）}
\begin{enumerate}[label=\Alph*., leftmargin=2.4em, itemsep=0.2em]
\item 0AH
\item 111D
\item 110Q
\item 1111B
\end{enumerate}
\end{question}
\begin{answer}
B（111D）
\end{answer}
\begin{explanation}
0AH=10，111D=111，110Q 是八进制 72，1111B 是二进制 15，因此最大的是 111D。
\end{explanation}

\begin{question}{6. 将超声波（振动波）转换成电信号是利用（）}
\begin{enumerate}[label=\Alph*., leftmargin=2.4em, itemsep=0.2em]
\item 电涡流效应
\item 压电效应
\item 应变效应
\item 逆压电效应
\end{enumerate}
\end{question}
\begin{answer}
B（压电效应）
\end{answer}
\begin{explanation}
压电效应指某些材料受机械振动或压力作用时产生电荷，可把超声振动转换为电信号。
\end{explanation}

\begin{question}{7. MPU6050 姿态传感器中，用于获取转动角速度的传感器单元是？}
\begin{enumerate}[label=\Alph*., leftmargin=2.4em, itemsep=0.2em]
\item I2C
\item DMP
\item Gyroscope
\item Accelerator
\end{enumerate}
\end{question}
\begin{answer}
C（Gyroscope）
\end{answer}
\begin{explanation}
Gyroscope 是陀螺仪，用来测量角速度；Accelerator/Accelerometer 测量加速度，DMP 是数字运动处理单元。
\end{explanation}

\begin{question}{8. 树莓派主板上使用的 SOC 主芯片内部的 CPU 采用的哪种架构？}
\begin{enumerate}[label=\Alph*., leftmargin=2.4em, itemsep=0.2em]
\item MIPS
\item RISCV
\item X86
\item ARM
\end{enumerate}
\end{question}
\begin{answer}
D（ARM）
\end{answer}
\begin{explanation}
树莓派采用 Broadcom SoC，CPU 架构属于 ARM 系列。
\end{explanation}

\begin{question}{9. 树莓派用来存储数据的媒介是}
\begin{enumerate}[label=\Alph*., leftmargin=2.4em, itemsep=0.2em]
\item Micro SD Card
\item 硬盘
\item SSD
\item 内存卡
\end{enumerate}
\end{question}
\begin{answer}
A（Micro SD Card）
\end{answer}
\begin{explanation}
树莓派通常通过 Micro SD Card 安装系统并存储数据。
\end{explanation}

\begin{question}{10. DHT11 温湿度传感器的数据接口是}
\begin{enumerate}[label=\Alph*., leftmargin=2.4em, itemsep=0.2em]
\item UART
\item SPI
\item 1-Wire
\item I2C
\end{enumerate}
\end{question}
\begin{answer}
C（1-Wire）
\end{answer}
\begin{explanation}
DHT11 使用单总线式数字通信，常被归为 1-Wire 类接口，而不是 UART、SPI 或 I2C。
\end{explanation}

\begin{question}{11. 下列在 OpenCV 中对色彩进行阈值化处理的函数是（）。}
\begin{enumerate}[label=\Alph*., leftmargin=2.4em, itemsep=0.2em]
\item cv2.GaussianBlur(src,ksize,sigmaX)
\item cv2.inRange(hsv,lower\_range,upper\_range)
\item cv2.cvtColor(input\_image,flag)
\item cv2.bitwise\_and(src1, src2, dst=None, mask=None)
\end{enumerate}
\end{question}
\begin{answer}
B（cv2.inRange(hsv,lower\_range,upper\_range)）
\end{answer}
\begin{explanation}
cv2.inRange() 用于判断像素是否落在给定颜色范围内，常用于 HSV 阈值分割。
\end{explanation}

\begin{question}{12. 以下内容关于函数描述正确的是？}
\begin{enumerate}[label=\Alph*., leftmargin=2.4em, itemsep=0.2em]
\item 以上说法都是正确的
\item 函数可以让程序执行的更快
\item 函数是一段代码用于执行特定的任务
\item 函数用于创建对象
\end{enumerate}
\end{question}
\begin{answer}
C（函数是一段代码用于执行特定的任务）
\end{answer}
\begin{explanation}
函数是封装一段完成特定任务的代码；函数本身不一定让程序更快，也不是专门用于创建对象。
\end{explanation}

\begin{question}{13. 假设要表示模拟输出有效，四路单端输入，禁止自动增量，A/D 通道为 0，那么该语 句的控制字是多少？bus.write\_byte(address,\_\_\_\_\_)}
\begin{enumerate}[label=\Alph*., leftmargin=2.4em, itemsep=0.2em]
\item 0x42
\item 0x41
\item 0x43
\item 0x40
\end{enumerate}
\end{question}
\begin{answer}
D（0x40）
\end{answer}
\begin{explanation}
PCF8591 控制字中 0x40 表示模拟输出使能、四路单端输入、禁止自动增量并选择 A/D 通道 0。
\end{explanation}

\begin{question}{14. PWM 波的高电平时间为 1ms，低电平时间为 3ms，它的占空比是多少？}
\begin{enumerate}[label=\Alph*., leftmargin=2.4em, itemsep=0.2em]
\item 无法确定
\item 75\%
\item 25\%
\item 33.30\%
\end{enumerate}
\end{question}
\begin{answer}
C（25\%）
\end{answer}
\begin{explanation}
周期 T=1ms+3ms=4ms，占空比=1/4=25%。
\end{explanation}

\begin{question}{15. 对于回流焊的过程的描述，以下哪一项是正确的。}
\begin{enumerate}[label=\Alph*., leftmargin=2.4em, itemsep=0.2em]
\item 焊锡膏经过预热 → 干燥 → 融化 → 润湿 → 冷却 → 流出回流焊炉
\item 焊锡膏经过预热 → 干燥 → 润湿 → 融化 → 冷却 → 流出回流焊炉
\item 焊锡膏经过干燥 → 融化 → 预热 → 润湿 → 冷却 → 流出回流焊炉
\item 焊锡膏经过干燥 → 预热 → 融化 → 润湿 → 冷却 → 流出回流焊炉
\end{enumerate}
\end{question}
\begin{answer}
D（焊锡膏经过干燥 → 预热 → 融化 → 润湿 → 冷却 → 流出回流焊炉）
\end{answer}
\begin{explanation}
回流焊通常经历干燥、预热、焊料融化、润湿、冷却等阶段，选项 D 顺序最符合过程。
\end{explanation}

\begin{question}{16. 在树莓派传感器实验中，很多输出开关量（只有高电平 1 或者低电平 0）输出的模块 中，常用的信号整形器件是（）。}
\begin{enumerate}[label=\Alph*., leftmargin=2.4em, itemsep=0.2em]
\item 三极管
\item 电阻
\item 二极管
\item 电压比较器
\end{enumerate}
\end{question}
\begin{answer}
D（电压比较器）
\end{answer}
\begin{explanation}
电压比较器能把模拟变化量整形成稳定的高/低电平开关量，适合数字输入模块。
\end{explanation}

\begin{question}{17. 电容的单位和符号分别是（）。}
\begin{enumerate}[label=\Alph*., leftmargin=2.4em, itemsep=0.2em]
\item 安培 A
\item 欧姆 Ω
\item 法拉 F
\item 伏特 V
\end{enumerate}
\end{question}
\begin{answer}
C（法拉 F）
\end{answer}
\begin{explanation}
电容的国际单位是法拉，符号为 F；安培是电流，欧姆是电阻，伏特是电压。
\end{explanation}

\begin{question}{18. 手工焊接时，焊锡在凝固之前，一经触动会产生（）。}
\begin{enumerate}[label=\Alph*., leftmargin=2.4em, itemsep=0.2em]
\item 虚焊点
\item 桥接
\item 短路
\item 完美焊点
\end{enumerate}
\end{question}
\begin{answer}
A（虚焊点）
\end{answer}
\begin{explanation}
焊锡未凝固前被触动会破坏焊点结晶和润湿状态，容易形成虚焊。
\end{explanation}

\begin{question}{19. 树莓派主板对外能输出的电源有（）}
\begin{enumerate}[label=\Alph*., leftmargin=2.4em, itemsep=0.2em]
\item 220V
\item 10V
\item 3.3V
\item 15V
\end{enumerate}
\end{question}
\begin{answer}
C（3.3V）
\end{answer}
\begin{explanation}
树莓派 GPIO 排针可提供 3.3V 和 5V 电源；在本题给出的选项中只有 3.3V 正确。
\end{explanation}

\begin{question}{20. 如图 PCF8591，当用跳线帽连接时，J4、J5、J6 分别对应的是哪三个传感器？}
\begin{enumerate}[label=\Alph*., leftmargin=2.4em, itemsep=0.2em]
\item 热敏电阻，电位计，光敏电阻
\item 光敏电阻，热敏电阻，电位计
\item 光敏电阻，电位计，热敏电阻
\item 电位计，光敏电阻，热敏电阻
\end{enumerate}
\end{question}
\begin{answer}
B（光敏电阻，热敏电阻，电位计）
\end{answer}
\begin{explanation}
该 PCF8591 实验板上 J4、J5、J6 跳线对应的模块顺序为光敏电阻、热敏电阻、电位计。
\end{explanation}

\begin{question}{21. PWM 波的高电平时间为 1ms，低电平时间为 3ms，它的频率是多少？图 1:}
\begin{enumerate}[label=\Alph*., leftmargin=2.4em, itemsep=0.2em]
\item 330Hz
\item 500Hz
\item 1KHz
\item 250Hz
\end{enumerate}
\end{question}
\begin{answer}
D（250Hz）
\end{answer}
\begin{explanation}
周期 T=1ms+3ms=4ms，频率 f=1/T=1/0.004s=250Hz。
\end{explanation}

\begin{question}{22. 树莓派对应的 TTL 电平中，高电平（1）对应的电压是多少？}
\begin{enumerate}[label=\Alph*., leftmargin=2.4em, itemsep=0.2em]
\item 2.0V
\item 15V
\item 0V
\item 3.3V
\end{enumerate}
\end{question}
\begin{answer}
D（3.3V）
\end{answer}
\begin{explanation}
树莓派 GPIO 逻辑高电平是 3.3V 级别，不是 5V 或更高电压。
\end{explanation}

\begin{question}{23. 当用跳线帽连接时，INPUT0、INPUT1、INPUT2 分别对应的是哪三个传感器？}
\begin{enumerate}[label=\Alph*., leftmargin=2.4em, itemsep=0.2em]
\item 光敏电阻，热敏电阻，电位计
\item 热敏电阻，电位计，光敏电阻
\item 电位计，光敏电阻，热敏电阻
\item 光敏电阻，电位计，热敏电阻
\end{enumerate}
\end{question}
\begin{answer}
C（电位计，光敏电阻，热敏电阻）
\end{answer}
\begin{explanation}
该模块跳线连接后 INPUT0、INPUT1、INPUT2 对应电位计、光敏电阻、热敏电阻。
\end{explanation}

\begin{question}{24. 关于 I2C 串行通信 SDA 和 SCL 的电路连接，下列说法错误的是？}
\begin{enumerate}[label=\Alph*., leftmargin=2.4em, itemsep=0.2em]
\item SCL 传送时钟信号
\item SDA 传送数据信号
\item SDA 和 SCL 数据线均通过上拉电阻连接到 VCC
\item SDA 和 SCL 数据线均通过下拉电阻连接到 VCC
\end{enumerate}
\end{question}
\begin{answer}
D（SDA 和 SCL 数据线均通过下拉电阻连接到 VCC）
\end{answer}
\begin{explanation}
I2C 的 SDA、SCL 通常通过上拉电阻接到 VCC；说成下拉电阻连接到 VCC 是错误的。
\end{explanation}

\begin{question}{25. I2C 串行通信总线空闲时，SDA 和 SCL 的电平状况分别是？}
\begin{enumerate}[label=\Alph*., leftmargin=2.4em, itemsep=0.2em]
\item SDA 为高电平，SCL 为低电平
\item SDA 为低电平，SCL 为低电平
\item SDA 为高电平，SCL 为高电平
\item SDA 为低电平，SCL 为高电平
\end{enumerate}
\end{question}
\begin{answer}
C（SDA 为高电平，SCL 为高电平）
\end{answer}
\begin{explanation}
I2C 总线空闲时，开漏结构由上拉电阻把 SDA 和 SCL 都拉为高电平。
\end{explanation}

\begin{question}{26. 关于 I2C 串行通信，下列说法不正确的是？}
\begin{enumerate}[label=\Alph*., leftmargin=2.4em, itemsep=0.2em]
\item 启动数据传送由主设备发起
\item I2C 通信中从设备都有一个唯一的地址
\item I2C 总线上的器件有主设备和从设备之分
\item 终止数据传送由从设备发起
\end{enumerate}
\end{question}
\begin{answer}
D（终止数据传送由从设备发起）
\end{answer}
\begin{explanation}
I2C 的起始和终止条件都由主设备控制发起，不是由从设备发起。
\end{explanation}

\begin{question}{27. 树莓派设置引脚为上拉输入模式时，相应的参数是？}
\begin{enumerate}[label=\Alph*., leftmargin=2.4em, itemsep=0.2em]
\item GPIO.OUT，GPIO.PUD\_UP
\item GPIO.OUT，GPIO.PUD\_DOWN
\item GPIO.IN，GPIO.PUD\_UP
\item GPIO.IN，GPIO.PUD\_DOWN
\end{enumerate}
\end{question}
\begin{answer}
C（GPIO.IN，GPIO.PUD\_UP）
\end{answer}
\begin{explanation}
上拉输入需要把引脚设置为输入模式 GPIO.IN，并启用 GPIO.PUD_UP。
\end{explanation}

\begin{question}{28. I2C 串行通信，用于传送数据的是？}
\begin{enumerate}[label=\Alph*., leftmargin=2.4em, itemsep=0.2em]
\item MOSI
\item SCL
\item SDA
\item MISO
\end{enumerate}
\end{question}
\begin{answer}
C（SDA）
\end{answer}
\begin{explanation}
I2C 中 SDA 是数据线，SCL 是时钟线；MOSI/MISO 属于 SPI。
\end{explanation}

\begin{question}{29. I2C 通信中，起始位是？}
\begin{enumerate}[label=\Alph*., leftmargin=2.4em, itemsep=0.2em]
\item SCL 为高电平，SDA 为上升沿
\item SCL 为高电平，SDA 为下降沿
\item SCL 为低电平，SDA 为下降沿
\item SCL 为低电平，SDA 为上升沿
\end{enumerate}
\end{question}
\begin{answer}
B（SCL 为高电平，SDA 为下降沿）
\end{answer}
\begin{explanation}
I2C 起始条件是在 SCL 为高电平时，SDA 从高电平跳变到低电平。
\end{explanation}

\begin{question}{30. 关于 Python 语言的特点，以下选项描述正确的是：}
\begin{enumerate}[label=\Alph*., leftmargin=2.4em, itemsep=0.2em]
\item Python 语言是解释型语言
\item Python 语言是编译型语言
\item Python 语言不支持面向对象
\item Python 语言是非跨平台语言
\end{enumerate}
\end{question}
\begin{answer}
A（Python 语言是解释型语言）
\end{answer}
\begin{explanation}
Python 主要采用解释执行方式，具有跨平台、支持面向对象等特点。
\end{explanation}

\begin{question}{31. 以下哪个不是 Python 的内置函数？}
\begin{enumerate}[label=\Alph*., leftmargin=2.4em, itemsep=0.2em]
\item char
\item float
\item int
\item list
\end{enumerate}
\end{question}
\begin{answer}
A（char）
\end{answer}
\begin{explanation}
int、float、list 都是 Python 内置函数/类型构造器，char 不是 Python 内置函数。
\end{explanation}

\begin{question}{32. 以下哪个属于 Python 合法的标识符？}
\begin{enumerate}[label=\Alph*., leftmargin=2.4em, itemsep=0.2em]
\item a@2
\item \_apple
\item global
\item 1\_ball
\end{enumerate}
\end{question}
\begin{answer}
B（\_apple）
\end{answer}
\begin{explanation}
标识符可以由字母、数字、下划线组成且不能以数字开头，也不能是关键字；_apple 合法。
\end{explanation}

\begin{question}{33. (2*5+3//2+15) and 10 的运算结果是多少？}
\begin{enumerate}[label=\Alph*., leftmargin=2.4em, itemsep=0.2em]
\item False
\item 20
\item 10
\item True
\end{enumerate}
\end{question}
\begin{answer}
C（10）
\end{answer}
\begin{explanation}
先算括号内为 2*5+3//2+15=26；在 Python 中 26 and 10 返回最后一个真值操作数 10。
\end{explanation}

\begin{question}{34. 如果函数没有使用 return 语句，则函数返回的是？}
\begin{enumerate}[label=\Alph*., leftmargin=2.4em, itemsep=0.2em]
\item 错误！函数必须要有返回值
\item 0
\item 任意的整数
\item None 对象
\end{enumerate}
\end{question}
\begin{answer}
D（None 对象）
\end{answer}
\begin{explanation}
Python 函数没有执行 return 或 return 后没有值时，默认返回 None。
\end{explanation}

\begin{question}{35. a=\{\}，判断 a 属于以下哪种数据类型？}
\begin{enumerate}[label=\Alph*., leftmargin=2.4em, itemsep=0.2em]
\item list
\item dict
\item set
\item tuple
\end{enumerate}
\end{question}
\begin{answer}
B（dict）
\end{answer}
\begin{explanation}
{} 表示空字典；空集合要写成 set()。
\end{explanation}

\begin{question}{36. 执行以下代码，依次输出哪些数字？f o r i i n range(2): p r i n t (i ) f o r i i n range(4,6): p r i n t (i )}
\begin{enumerate}[label=\Alph*., leftmargin=2.4em, itemsep=0.2em]
\item 0,1,2,4,5,6
\item 2,4,6
\item 1,2,4,5,6
\item 0,1,4,5
\end{enumerate}
\end{question}
\begin{answer}
D（0,1,4,5）
\end{answer}
\begin{explanation}
range(2) 产生 0、1；range(4,6) 产生 4、5，因此依次输出 0、1、4、5。
\end{explanation}

\begin{question}{37. Python 用以下哪个关键字来声明一个类？}
\begin{enumerate}[label=\Alph*., leftmargin=2.4em, itemsep=0.2em]
\item class
\item block
\item def
\item object
\end{enumerate}
\end{question}
\begin{answer}
A（class）
\end{answer}
\begin{explanation}
Python 使用 class 关键字定义类，def 用于定义函数。
\end{explanation}

\begin{question}{38. 电阻的单位是（）。}
\begin{enumerate}[label=\Alph*., leftmargin=2.4em, itemsep=0.2em]
\item 亨特
\item 欧姆
\item 安培
\item 法拉
\end{enumerate}
\end{question}
\begin{answer}
B（欧姆）
\end{answer}
\begin{explanation}
电阻的单位是欧姆，符号常写作 欧姆。
\end{explanation}

\begin{question}{39. 在润湿作用中，润湿角一般在（）范围内最佳。}
\begin{enumerate}[label=\Alph*., leftmargin=2.4em, itemsep=0.2em]
\item 45 ◦ –90 ◦
\item 大于 90 ◦
\item 15 ◦ –45 ◦
\item 0 ◦ –15 ◦
\end{enumerate}
\end{question}
\begin{answer}
C（15 ◦ –45 ◦ ）
\end{answer}
\begin{explanation}
润湿角过大说明润湿差，通常 15度-45度 被认为润湿效果较好。
\end{explanation}

\begin{question}{40. 清洁电烙铁头时，应该将电烙铁温度调至（）}
\begin{enumerate}[label=\Alph*., leftmargin=2.4em, itemsep=0.2em]
\item 350◦ C
\item 250◦ C
\item 400◦ C
\item 300◦ C
\end{enumerate}
\end{question}
\begin{answer}
B（250◦ C）
\end{answer}
\begin{explanation}
清洁烙铁头时温度不宜过高，约 250摄氏度 有利于减少氧化并进行清洁。
\end{explanation}

\begin{question}{41. AOI 光学检测的流程包括以下六项：1 缺陷处理；2 板件扫描；3 影像处理；4 资料 处理；5 母版学习；6 逻辑对比。其正确的流程顺序为：}
\begin{enumerate}[label=\Alph*., leftmargin=2.4em, itemsep=0.2em]
\item 4 → 5 → 2 → 3 → 6 → 1
\item 4 → 5 → 2 → 3 → 1 → 6
\item 5 → 4 → 2 → 3 → 6 → 1
\item 5 → 4 → 2 → 3 → 1 → 6
\end{enumerate}
\end{question}
\begin{answer}
A（4 → 5 → 2 → 3 → 6 → 1）
\end{answer}
\begin{explanation}
AOI 通常先进行资料处理和母版学习，再扫描板件、处理影像、逻辑对比，最后缺陷处理。
\end{explanation}

\begin{question}{42. 调节刮板角度范围一般为：}
\begin{enumerate}[label=\Alph*., leftmargin=2.4em, itemsep=0.2em]
\item 30 ◦ –40 ◦
\item 75 ◦ –90 ◦
\item 40 ◦ –75 ◦
\item 0 ◦ –30 ◦
\end{enumerate}
\end{question}
\begin{answer}
C（40 ◦ –75 ◦ ）
\end{answer}
\begin{explanation}
刮板角度一般调节在 40度-75度，角度会影响锡膏滚动和印刷质量。
\end{explanation}

\begin{question}{43. 树莓派数据上云实验中，树莓派和阿里云物联网平台之间数据传输协议是（）。}
\begin{enumerate}[label=\Alph*., leftmargin=2.4em, itemsep=0.2em]
\item MQTT
\item CoAP
\item XMPP
\item DDS
\end{enumerate}
\end{question}
\begin{answer}
A（MQTT）
\end{answer}
\begin{explanation}
阿里云物联网平台常用 MQTT 进行设备端和云端的发布/订阅式数据传输。
\end{explanation}

\begin{question}{44. 在 Linux 系统中，可以显示当前系统未使用的和已使用的内存数目的命令是（）。}
\begin{enumerate}[label=\Alph*., leftmargin=2.4em, itemsep=0.2em]
\item cat
\item free
\item df
\item top
\end{enumerate}
\end{question}
\begin{answer}
B（free）
\end{answer}
\begin{explanation}
free 命令用于查看系统内存使用与空闲情况；df 查看磁盘，top 查看进程资源。
\end{explanation}

\begin{question}{45. imutils.translate(img, x, y) 表示（）。}
\begin{enumerate}[label=\Alph*., leftmargin=2.4em, itemsep=0.2em]
\item 图像缩放
\item 图像旋转
\item 图像平移
\item 骨架提取
\end{enumerate}
\end{question}
\begin{answer}
C（图像平移）
\end{answer}
\begin{explanation}
translate 表示平移，imutils.translate(img,x,y) 将图像在 x、y 方向移动。
\end{explanation}

\begin{question}{46. 舵机转速定义为舵机无负载情况下，转过（）度角所需时间。}
\begin{enumerate}[label=\Alph*., leftmargin=2.4em, itemsep=0.2em]
\item 0
\item 90
\item 180
\item 60
\end{enumerate}
\end{question}
\begin{answer}
D（60）
\end{answer}
\begin{explanation}
舵机速度常以无负载转过 60度 所需时间表示，例如 0.12s/60度。 二、多选题
\end{explanation}

\section{二、多选题}

\begin{question}{47. 在 Python 中，字符串 var1 = 'abcde'，想要从 var1 中截取字符串"bcd"，以下 表达正确的是？}
\begin{enumerate}[label=\Alph*., leftmargin=2.4em, itemsep=0.2em]
\item var1[1:-1]
\item var1[-4:-1]
\item var1[2:4]
\item var1[1:4]
\end{enumerate}
\end{question}
\begin{answer}
A（var1[1:-1]），B（var1[-4:-1]），D（var1[1:4]）
\end{answer}
\begin{explanation}
字符串切片左闭右开；abcde 中 bcd 的下标范围是 1 到 4，也可用负下标 -4 到 -1。
\end{explanation}

\begin{question}{48. 树莓派的编码方式有}
\begin{enumerate}[label=\Alph*., leftmargin=2.4em, itemsep=0.2em]
\item BCM 编码
\item WiringPi 编码
\item 以上都不是
\item BOARD 编码
\end{enumerate}
\end{question}
\begin{answer}
A（BCM 编码），B（WiringPi 编码），D（BOARD 编码）
\end{answer}
\begin{explanation}
树莓派 GPIO 常见编号方式包括 BOARD 物理引脚编号、BCM 编号和 WiringPi 编号。
\end{explanation}

\begin{question}{49. 在做三色 LED 实验时，需要（）彩色分量，通过 RGB 三个分量的不同比例的组合， 可得到任意的颜色。}
\begin{enumerate}[label=\Alph*., leftmargin=2.4em, itemsep=0.2em]
\item 蓝
\item 绿
\item 黄
\item 红
\end{enumerate}
\end{question}
\begin{answer}
A（蓝），B（绿），D（红）
\end{answer}
\begin{explanation}
RGB 三色 LED 的三个基本颜色分量是红、绿、蓝，按比例混合可形成多种颜色。
\end{explanation}

\begin{question}{50. MPU6050 得到的姿态角包含？}
\begin{enumerate}[label=\Alph*., leftmargin=2.4em, itemsep=0.2em]
\item Pitch
\item gyroscope
\item yaw
\item roll
\end{enumerate}
\end{question}
\begin{answer}
A（Pitch），C（yaw），D（roll）
\end{answer}
\begin{explanation}
姿态角通常包括 pitch、roll、yaw；gyroscope 是传感器单元，不是姿态角名称。
\end{explanation}

\begin{question}{51. 电感式传感器可以对（）等物理量进行测量。}
\begin{enumerate}[label=\Alph*., leftmargin=2.4em, itemsep=0.2em]
\item 振动
\item 位移
\item 压力
\item 流量
\end{enumerate}
\end{question}
\begin{answer}
A（振动），B（位移），C（压力） ，D（流量）
\end{answer}
\begin{explanation}
电感式传感器可把位移等变化转换成电感变化，进一步测量振动、压力、流量等物理量。
\end{explanation}

\begin{question}{52. 关于 I2C 总线数据传输，下列说法正确的是（）}
\begin{enumerate}[label=\Alph*., leftmargin=2.4em, itemsep=0.2em]
\item SCL 为高电平期间，SDA 由高电平向低电平的变化表示起始信号
\item SCL 为高电平期间，SDA 由高电平向低电平的变化表示终止信号
\item SCL 为高电平期间，SDA 由低电平向高电平的变化表示终止信号
\item SCL 为高电平期间，SDA 由低电平向高电平的变化表示起始信号
\end{enumerate}
\end{question}
\begin{answer}
A（SCL 为高电平期间，SDA 由高电平向低电平的变化表示起始信号） ，C（SCL 为 高电平期间，SDA 由低电平向高电平的变化表示终止信号）
\end{answer}
\begin{explanation}
I2C 中 SCL 高电平时 SDA 下降沿表示起始信号，上升沿表示终止信号。
\end{explanation}

\begin{question}{53. 影响焊接界面结合层的因素有：}
\begin{enumerate}[label=\Alph*., leftmargin=2.4em, itemsep=0.2em]
\item 助焊剂的质量
\item 焊料的合金成分和氧化程度
\item 焊接的温度和时间
\item 母材的氧化程度
\end{enumerate}
\end{question}
\begin{answer}
A（助焊剂的质量） ，C（焊接的温度和时间） ，B（焊料的合金成分和氧化程度） ，D （母材的氧化程度）
\end{answer}
\begin{explanation}
焊接界面结合层受助焊剂、焊料成分与氧化程度、温度时间、母材氧化等多因素影响。
\end{explanation}

\begin{question}{54. 在不增加硬件的情况下，树莓派本体能提供的 IO 接口的功能类型有}
\begin{enumerate}[label=\Alph*., leftmargin=2.4em, itemsep=0.2em]
\item CAN
\item GPIO
\item I2C
\item SPI
\end{enumerate}
\end{question}
\begin{answer}
B（GPIO），C（I2C），D（SPI）
\end{answer}
\begin{explanation}
树莓派本体常用接口包括 GPIO、I2C、SPI；CAN 不是其默认直接提供的 40Pin 常规接口。
\end{explanation}

\begin{question}{55. PCF8591 具备功能接口有}
\begin{enumerate}[label=\Alph*., leftmargin=2.4em, itemsep=0.2em]
\item AD 转换
\item GPIO
\item DA 转换
\item I2C 接口
\end{enumerate}
\end{question}
\begin{answer}
A（AD 转换），C（DA 转换） ，D（I2C 接口）
\end{answer}
\begin{explanation}
PCF8591 集成 A/D、D/A 以及 I2C 接口，不是通用 GPIO 扩展芯片。
\end{explanation}

\begin{question}{56. 树莓派中 GPIO 口具备 PWM 功能，如果想使用 PWM 功能需要设置的参数有}
\begin{enumerate}[label=\Alph*., leftmargin=2.4em, itemsep=0.2em]
\item PWM 的频率
\item 对应 IO 要设置成输入模式
\item PWM 的占空比
\item 对应 IO 要设置成输出模式
\end{enumerate}
\end{question}
\begin{answer}
A（PWM 的频率），C（PWM 的占空比），D（对应 IO 要设置成输出模式）
\end{answer}
\begin{explanation}
使用 PWM 需要选择输出引脚、设置 PWM 频率并调整占空比；输入模式不能输出 PWM。
\end{explanation}

\begin{question}{57. 树莓派的板载接口包括（）。}
\begin{enumerate}[label=\Alph*., leftmargin=2.4em, itemsep=0.2em]
\item 1-Wire 接口
\item I2C 接口
\item SPI 接口
\item 通用 IO 口
\end{enumerate}
\end{question}
\begin{answer}
A（1-Wire 接口），B（I2C 接口），C（SPI 接口），D（通用 IO 口）
\end{answer}
\begin{explanation}
树莓派 40Pin 可复用 GPIO、I2C、SPI、1-Wire 等接口功能。
\end{explanation}

\begin{question}{58. 关于 return 说法正确的是？}
\begin{enumerate}[label=\Alph*., leftmargin=2.4em, itemsep=0.2em]
\item return 没有返回值时，函数自动返回 None。
\item Python 函数中必须有 return。
\item 执行到 return 时，程序将停止执行函数内 return 后边的语句。
\item return 可以返回多个值。
\end{enumerate}
\end{question}
\begin{answer}
A（return 没有返回值时，函数自动返回 None。 ），C（执行到 return 时，程序将停 止执行函数内 return 后边的语句。），D（return 可以返回多个值。）
\end{answer}
\begin{explanation}
return 可返回值或多个值；执行 return 后函数结束；没有返回值时结果为 None。
\end{explanation}

\begin{question}{59. 以下哪些属于正确的 Python 注释方式？}
\begin{enumerate}[label=\Alph*., leftmargin=2.4em, itemsep=0.2em]
\item \%comments
\item "'comments"'
\item //comments
\item \#comments
\end{enumerate}
\end{question}
\begin{answer}
B（"'comments"'），D（\#comments）
\end{answer}
\begin{explanation}
Python 单行注释用 #；三引号字符串常用于多行注释式说明，// 和 % 不是 Python 注释符。
\end{explanation}

\begin{question}{60. PWM 对应的参数有哪些？}
\begin{enumerate}[label=\Alph*., leftmargin=2.4em, itemsep=0.2em]
\item 高电平
\item 周期
\item 频率
\item 占空比
\end{enumerate}
\end{question}
\begin{answer}
B（周期），C（频率），D（占空比）
\end{answer}
\begin{explanation}
PWM 的核心参数包括周期、频率和占空比；高电平只是波形状态，不是通常说的独立参数。
\end{explanation}

\begin{question}{61. AOI 光学检测的构成包括以下哪几部分？}
\begin{enumerate}[label=\Alph*., leftmargin=2.4em, itemsep=0.2em]
\item 软件系统——分析处理
\item 视觉系统——图像采集
\item 机械系统——传送
\item 刮板——刮动焊锡膏
\end{enumerate}
\end{question}
\begin{answer}
A（软件系统——分析处理） ，C（机械系统——传送） ，B（视觉系统——图像采集）
\end{answer}
\begin{explanation}
AOI 系统主要由视觉采集、机械传送和软件分析处理构成；刮板属于印刷设备部件。
\end{explanation}

\begin{question}{62. 在使用树莓派的过程中，树莓派的 40Pin 的接口相关的编址编码有（）。}
\begin{enumerate}[label=\Alph*., leftmargin=2.4em, itemsep=0.2em]
\item BCM
\item BOARD
\item ERP
\item CRM
\end{enumerate}
\end{question}
\begin{answer}
A（BCM），B（BOARD）
\end{answer}
\begin{explanation}
树莓派 40Pin 常见编址方式是 BCM 编号和 BOARD 物理编号。
\end{explanation}

\begin{question}{63. 二极管的作用包括以下哪几种？}
\begin{enumerate}[label=\Alph*., leftmargin=2.4em, itemsep=0.2em]
\item 开关
\item 限幅
\item 整流
\item 隔直通交
\end{enumerate}
\end{question}
\begin{answer}
A（开关），B（限幅），C（整流）
\end{answer}
\begin{explanation}
二极管具有单向导电性，可用于整流、限幅、开关等；隔直通交通常是电容的作用。
\end{explanation}

\begin{question}{64. 如下通信接口中是同步串行通信的是（）}
\begin{enumerate}[label=\Alph*., leftmargin=2.4em, itemsep=0.2em]
\item SPI
\item I2C
\item 1-Wire
\item UART
\end{enumerate}
\end{question}
\begin{answer}
A（SPI），B（I2C）
\end{answer}
\begin{explanation}
SPI 和 I2C 都有时钟线，属于同步串行通信；UART 是异步串行通信。
\end{explanation}

\begin{question}{65. 传感器由（）组成。}
\begin{enumerate}[label=\Alph*., leftmargin=2.4em, itemsep=0.2em]
\item 测量电路
\item 反馈电路
\item 敏感元件
\item 转换元件
\end{enumerate}
\end{question}
\begin{answer}
A（测量电路），C（敏感元件） ，D（转换元件）
\end{answer}
\begin{explanation}
传感器一般由敏感元件、转换元件和测量电路组成。
\end{explanation}

\begin{question}{66. 树莓派的管脚有哪些功能可以使用？}
\begin{enumerate}[label=\Alph*., leftmargin=2.4em, itemsep=0.2em]
\item 通用 IO （GPIO）
\item I2C 口
\item 串口
\item 1-Wire 接口（单总线）
\end{enumerate}
\end{question}
\begin{answer}
A（通用 IO （GPIO）），B（I2C 口），C（串口），D（1-Wire 接口（单总线））
\end{answer}
\begin{explanation}
树莓派管脚可复用为 GPIO、I2C、串口、1-Wire 等功能。
\end{explanation}

\begin{question}{67. 树莓派的管脚有哪些功能可以使用？}
\begin{enumerate}[label=\Alph*., leftmargin=2.4em, itemsep=0.2em]
\item 1-Wire 接口（单总线）
\item SPI 口
\item I2C 口
\item 串口
\end{enumerate}
\end{question}
\begin{answer}
A（1-Wire 接口（单总线）），B（SPI 口），C（I2C 口），D（串口）
\end{answer}
\begin{explanation}
树莓派管脚可复用为 1-Wire、SPI、I2C、串口等功能。
\end{explanation}

\begin{question}{68. 树莓派引脚设置中，用于设置引脚输出状态的有哪些？}
\begin{enumerate}[label=\Alph*., leftmargin=2.4em, itemsep=0.2em]
\item INPUT\_PULLDOWN
\item GPIO.LOW
\item GPIO.HIGH
\item INPUT\_PULLUP
\end{enumerate}
\end{question}
\begin{answer}
B（GPIO.LOW），C（GPIO.HIGH）
\end{answer}
\begin{explanation}
设置输出状态就是输出高电平或低电平，即 GPIO.HIGH 或 GPIO.LOW。
\end{explanation}

\begin{question}{69. GPIO.PWM() 的参数分别为？}
\begin{enumerate}[label=\Alph*., leftmargin=2.4em, itemsep=0.2em]
\item 占空比
\item 起始标志
\item PWM 频率
\item 引脚号
\end{enumerate}
\end{question}
\begin{answer}
C（PWM 频率），D（引脚号）
\end{answer}
\begin{explanation}
GPIO.PWM(channel, frequency) 的两个参数通常是引脚号和 PWM 频率。
\end{explanation}

\begin{question}{70. 在 Python 中，以下哪些属于可变的数据类型？}
\begin{enumerate}[label=\Alph*., leftmargin=2.4em, itemsep=0.2em]
\item list
\item tuple
\item set
\item dict
\end{enumerate}
\end{question}
\begin{answer}
A（list），C（set），D（dict）
\end{answer}
\begin{explanation}
list、set、dict 都是可变类型；tuple 是不可变类型。
\end{explanation}

\begin{question}{71. 电阻的作用包括以下哪些？}
\begin{enumerate}[label=\Alph*., leftmargin=2.4em, itemsep=0.2em]
\item 耦合
\item 分压
\item 限流
\item 滤波
\end{enumerate}
\end{question}
\begin{answer}
A（耦合），B（分压），C（限流） ，D（滤波）
\end{answer}
\begin{explanation}
电阻可用于限流、分压，也可与其他元件组成耦合、滤波等电路。
\end{explanation}

\begin{question}{72. 什么情况会造成电解电容爆裂？}
\begin{enumerate}[label=\Alph*., leftmargin=2.4em, itemsep=0.2em]
\item 正负反接
\item 超过耐压值
\item 正负正接
\item 不超过耐压值
\end{enumerate}
\end{question}
\begin{answer}
A（正负反接），B（超过耐压值）
\end{answer}
\begin{explanation}
电解电容有极性和耐压限制，反接或超过耐压都可能导致发热、漏液甚至爆裂。
\end{explanation}

\begin{question}{73. 刮板及参数的调节中，要调节的参数包括以下哪些？}
\begin{enumerate}[label=\Alph*., leftmargin=2.4em, itemsep=0.2em]
\item 加速度
\item 角度
\item 速度
\item 压力
\end{enumerate}
\end{question}
\begin{answer}
B（角度），C（速度），D（压力）
\end{answer}
\begin{explanation}
刮板调节通常关注角度、速度和压力，这些直接影响锡膏印刷厚度与均匀性。
\end{explanation}

\begin{question}{74. 焊接中，减小表面张力的方法有：}
\begin{enumerate}[label=\Alph*., leftmargin=2.4em, itemsep=0.2em]
\item 降低焊接温度
\item 减少活性剂
\item 增加活性剂
\item 提高焊接温度
\end{enumerate}
\end{question}
\begin{answer}
C（增加活性剂），D（提高焊接温度）
\end{answer}
\begin{explanation}
提高温度、增加活性剂能改善润湿并减小表面张力；降低温度或减少活性剂效果相反。
\end{explanation}

\begin{question}{75. 以下特点中，属于 SMT 工艺优点的有：}
\begin{enumerate}[label=\Alph*., leftmargin=2.4em, itemsep=0.2em]
\item 成本高
\item 高频特性好
\item 组装密度高
\item 可靠性高
\end{enumerate}
\end{question}
\begin{answer}
B（高频特性好），C（组装密度高），D（可靠性高）
\end{answer}
\begin{explanation}
SMT 具有组装密度高、高频特性好、可靠性高等优点；成本高不是优点。
\end{explanation}

\begin{question}{76. 以下属于物联网网络模型的是（）}
\begin{enumerate}[label=\Alph*., leftmargin=2.4em, itemsep=0.2em]
\item 感知层
\item 网络层
\item 平台层
\item 应用层
\end{enumerate}
\end{question}
\begin{answer}
A（感知层），B（网络层），C（平台层），D（应用层）
\end{answer}
\begin{explanation}
物联网常分为感知层、网络层、平台层和应用层。
\end{explanation}

\begin{question}{77. 在树莓派实验中，阿里云设备三元组信息包括（）。}
\begin{enumerate}[label=\Alph*., leftmargin=2.4em, itemsep=0.2em]
\item ProductKey
\item DeviceName
\item DeviceSecret
\item ProductSecret
\end{enumerate}
\end{question}
\begin{answer}
A（ProductKey），B（DeviceName），C（DeviceSecret）
\end{answer}
\begin{explanation}
阿里云设备三元组是 ProductKey、DeviceName、DeviceSecret。
\end{explanation}

\begin{question}{78. 使用 argparse 的主要步骤是（）。}
\begin{enumerate}[label=\Alph*., leftmargin=2.4em, itemsep=0.2em]
\item 导入 argparse 包
\item 创建 ArgumentParser() 参数对象
\item 调用 add\_argument() 方法往参数对象中添加参数
\item 使用 parse\_args() 解析添加参数的参数对象，获得解析对象
\end{enumerate}
\end{question}
\begin{answer}
A（导入 argparse 包） ，C（调用 add\_argument() ，B（创建 ArgumentParser() 参数对象） ，D（使用 parse\_args() 解析添加参数的参数对象，获得解析 方法往参数对象中添加参数） 对象）
\end{answer}
\begin{explanation}
argparse 的使用流程是导入包、创建解析器、添加参数、parse_args() 解析。
\end{explanation}

\begin{question}{79. 舵机驱动板 PCA9685 的主要特点是（）}
\begin{enumerate}[label=\Alph*., leftmargin=2.4em, itemsep=0.2em]
\item I2C 接口，支持高达 16 路 PWM 输出，每路 12 位分辨率（4096 级）
\item 内置 25MHz 晶振，也可使用外部晶振，最大 50MHz
\item 支持 2.3V-5.5V 电压，最大耐压值 5.5V，逻辑电平 3.3V
\item 具有上电复位，以及软件复位等功能
\end{enumerate}
\end{question}
\begin{answer}
A（I2C 接口，支持高达 16 路 PWM 输出，每路 12 位分辨率（4096 级） ），B（内 置 25MHz 晶振，也可使用外部晶振，最大 50MHz），C（支持 2.3V-5.5V 电压，最大耐压 值 5.5V，逻辑电平 3.3V），D（具有上电复位，以及软件复位等功能）
\end{answer}
\begin{explanation}
PCA9685 是 I2C 控制的 16 路 12 位 PWM 驱动芯片，支持内部/外部时钟、低压逻辑和复位功能。 三、判断题
\end{explanation}

\section{三、判断题}

\begin{question}{80. I2C 通信中，SDA 数据线仅用于发送数据。}
\begin{enumerate}[label=\Alph*., leftmargin=2.4em, itemsep=0.2em]
\item 正确
\item 错误
\end{enumerate}
\end{question}
\begin{answer}
B（错误） 代码可写为 GPIO.setmode(GPIO.BOARD)。
\end{answer}
\begin{explanation}
SDA 是双向数据线，可发送也可接收数据，因此“仅用于发送”错误。
\end{explanation}

\begin{question}{81. 树莓派 GPIO 设置为 Board 编码时，代码可写为 GPIO.setmode(GPIO.BOARD)。}
\begin{enumerate}[label=\Alph*., leftmargin=2.4em, itemsep=0.2em]
\item 正确
\item 错误
\end{enumerate}
\end{question}
\begin{answer}
A（正确）
\end{answer}
\begin{explanation}
使用 BOARD 物理编号时，代码写 GPIO.setmode(GPIO.BOARD)。
\end{explanation}

\begin{question}{82. PCF8591 的 A0、A1 和 A2 用于硬件地址的编程，最多在同个 I2C 总线上可接入 8 个 PCF8591。}
\begin{enumerate}[label=\Alph*., leftmargin=2.4em, itemsep=0.2em]
\item 正确
\item 错误
\end{enumerate}
\end{question}
\begin{answer}
A（正确）
\end{answer}
\begin{explanation}
A0、A1、A2 三个硬件地址位可组合 2^3=8 个地址，所以同一 I2C 总线最多可区分 8 个 PCF8591。
\end{explanation}

\begin{question}{83. 热敏电阻除了用于温度传感器以外，也可用于湿度传感器。}
\begin{enumerate}[label=\Alph*., leftmargin=2.4em, itemsep=0.2em]
\item 正确
\item 错误
\end{enumerate}
\end{question}
\begin{answer}
A（正确）
\end{answer}
\begin{explanation}
热敏电阻主要用于温度检测，在部分教学题库中也可作为湿度检测相关电路的补偿或辅助元件； 本题按“正确”处理。
\end{explanation}

\begin{question}{84. SPI、UART、I2C 都是串行数据总线。}
\begin{enumerate}[label=\Alph*., leftmargin=2.4em, itemsep=0.2em]
\item 错误
\item 正确
\end{enumerate}
\end{question}
\begin{answer}
B（正确）
\end{answer}
\begin{explanation}
SPI、UART、I2C 都按位串行传输数据，因此都属于串行数据总线。
\end{explanation}

\begin{question}{85. LCD 液晶模块上的电位器是用来调整显示对比度的。}
\begin{enumerate}[label=\Alph*., leftmargin=2.4em, itemsep=0.2em]
\item 正确
\item 错误
\end{enumerate}
\end{question}
\begin{answer}
A（正确）
\end{answer}
\begin{explanation}
LCD1602 等液晶模块上的电位器通常用于调节液晶显示对比度。
\end{explanation}

\begin{question}{86. 多彩 LED 灯是由红、绿、蓝三个颜色的 LED 灯珠组合而成。}
\begin{enumerate}[label=\Alph*., leftmargin=2.4em, itemsep=0.2em]
\item 错误
\item 正确
\end{enumerate}
\end{question}
\begin{answer}
B（正确）
\end{answer}
\begin{explanation}
多彩/RGB LED 由红、绿、蓝三种 LED 芯片组合，通过混色产生不同颜色。
\end{explanation}

\begin{question}{87. 运行 import numpy as np 之后，我们在调用 numpy 的内置函数时，用 np. 加函 数名即可。}
\begin{enumerate}[label=\Alph*., leftmargin=2.4em, itemsep=0.2em]
\item 正确
\item 错误
\end{enumerate}
\end{question}
\begin{answer}
A（正确）
\end{answer}
\begin{explanation}
import numpy as np 给 numpy 起别名 np，调用函数时写 np.函数名。
\end{explanation}

\begin{question}{88. Python 表达式 list(range(5)) 的值为 [1,2,3,4,5]。}
\begin{enumerate}[label=\Alph*., leftmargin=2.4em, itemsep=0.2em]
\item 错误
\item 正确
\end{enumerate}
\end{question}
\begin{answer}
A（错误）
\end{answer}
\begin{explanation}
range(5) 产生 0 到 4，list(range(5)) 为 [0,1,2,3,4]，不是 [1,2,3,4,5]。
\end{explanation}

\begin{question}{89. 在 Python 中可以使用 if 作为变量名。}
\begin{enumerate}[label=\Alph*., leftmargin=2.4em, itemsep=0.2em]
\item 错误
\item 正确
\end{enumerate}
\end{question}
\begin{answer}
A（错误）
\end{answer}
\begin{explanation}
if 是 Python 保留关键字，不能作为变量名。
\end{explanation}

\begin{question}{90. 用频率为 2KHz 的 PWM 波控制 LED 时，看不出灯有闪烁，是利用了人眼的视觉 暂留特性。}
\begin{enumerate}[label=\Alph*., leftmargin=2.4em, itemsep=0.2em]
\item 错误
\item 正确
\end{enumerate}
\end{question}
\begin{answer}
B（正确）
\end{answer}
\begin{explanation}
2KHz 远高于人眼可分辨闪烁频率，LED 看起来连续发光，是视觉暂留作用。
\end{explanation}

\begin{question}{91. 在 try...except...else 结构中，如果 try 块的语句引发了异常，则会执行 else 块中的 语句。}
\begin{enumerate}[label=\Alph*., leftmargin=2.4em, itemsep=0.2em]
\item 错误
\item 正确
\end{enumerate}
\end{question}
\begin{answer}
A（错误）
\end{answer}
\begin{explanation}
try...except...else 中，else 只在 try 没有异常时执行；发生异常则进入 except。
\end{explanation}

\begin{question}{92. 传感器是一种把非电量的被测量转换为与之确定对应关系的、便于应用的电量的测量 装置。}
\begin{enumerate}[label=\Alph*., leftmargin=2.4em, itemsep=0.2em]
\item 正确
\item 错误
\end{enumerate}
\end{question}
\begin{answer}
A（正确）
\end{answer}
\begin{explanation}
传感器的作用就是把被测非电量转换为便于处理的电信号或相应输出量。
\end{explanation}

\begin{question}{93. GPIO 口是通用输入输出的简称，PIN 脚依照考量可作为通用输入（GPI）或通用输 出（GPO）或者通用输入输出（GPIO）。}
\begin{enumerate}[label=\Alph*., leftmargin=2.4em, itemsep=0.2em]
\item 正确
\item 错误
\end{enumerate}
\end{question}
\begin{answer}
A（正确）
\end{answer}
\begin{explanation}
GPIO 是 General Purpose Input/Output，可配置为输入、输出或具备复用功能。
\end{explanation}

\begin{question}{94. 舵机是集成了直流电机、电机控制器和减速器等，并封装在一个便于安装的外壳里的 伺服单元，能够利用简单的输入信号比较精确地转动给定角度的电机系统。}
\begin{enumerate}[label=\Alph*., leftmargin=2.4em, itemsep=0.2em]
\item 正确
\item 错误
\end{enumerate}
\end{question}
\begin{answer}
A（正确）
\end{answer}
\begin{explanation}
舵机集成电机、减速机构和控制电路，可根据控制信号转到指定角度。
\end{explanation}

\begin{question}{95. continue 语句的作用是提前结束本次循环，忽略之后的所有语句，提前进入下一次循 环。}
\begin{enumerate}[label=\Alph*., leftmargin=2.4em, itemsep=0.2em]
\item 正确
\item 错误
\end{enumerate}
\end{question}
\begin{answer}
A（正确） 低功耗、
\end{answer}
\begin{explanation}
continue 会跳过本次循环剩余语句，直接进入下一轮循环判断。
\end{explanation}

\begin{question}{96. PCF8591 是一个单片集成、单独供电、 12 位 CMOS 数据获取器件。 PCF8591 其中具有 4 个模拟输入、1 个模拟输出和 1 个串行 I2C 总线接口}
\begin{enumerate}[label=\Alph*., leftmargin=2.4em, itemsep=0.2em]
\item 正确
\item 错误
\end{enumerate}
\end{question}
\begin{answer}
B（错误）
\end{answer}
\begin{explanation}
PCF8591 是 8 位 CMOS 数据采集器件，不是 12 位；它有 4 路模拟输入、1 路模拟输出和 I2C 接口。
\end{explanation}

\begin{question}{97. imutils 是 Adrian Rosebrock 开发的一个 python 工具包，它整合了 opencv、 numpy 和 matplotlib 的相关操作，主要是用来进行图形图像的处理，后期又加入了针对 视频的处理等。}
\begin{enumerate}[label=\Alph*., leftmargin=2.4em, itemsep=0.2em]
\item 正确
\item 错误
\end{enumerate}
\end{question}
\begin{answer}
A（正确）
\end{answer}
\begin{explanation}
imutils 是常用 OpenCV 辅助工具包，封装了图像平移、旋转、缩放及视频处理等便捷操作。
\end{explanation}

\begin{question}{98. rotated\_round = imutils.rotate\_bound(image, 90) 表示图像逆时针旋转 90 度。}
\begin{enumerate}[label=\Alph*., leftmargin=2.4em, itemsep=0.2em]
\item 正确
\item 错误
\end{enumerate}
\end{question}
\begin{answer}
B（错误）
\end{answer}
\begin{explanation}
imutils.rotate_bound(image,90) 在常用坐标约定下表示将图像逆时针旋转 90度 并自动扩展边界。
\end{explanation}

\begin{question}{99. 只要二极管的正极电压比负极电压高，二极管就一定导通。}
\begin{enumerate}[label=\Alph*., leftmargin=2.4em, itemsep=0.2em]
\item 错误
\item 正确
\end{enumerate}
\end{question}
\begin{answer}
A（错误）
\end{answer}
\begin{explanation}
二极管导通还需要正向电压达到导通压降，并受器件材料和电路条件限制，不是只要正极略高就 一定导通。
\end{explanation}

\begin{question}{100. 在电路搭建好之后，确认无误，才可以给电路供电。}
\begin{enumerate}[label=\Alph*., leftmargin=2.4em, itemsep=0.2em]
\item 正确
\item 错误
\end{enumerate}
\end{question}
\begin{answer}
A（正确）
\end{answer}
\begin{explanation}
电路搭建完成后必须先检查连接、极性和电源电压，确认无误后再上电，避免短路或烧毁元件。
\end{explanation}

\begin{question}{101. 仅通过 MPU6050 姿态传感器获得的角速度数据可以解析出姿态角。}
\begin{enumerate}[label=\Alph*., leftmargin=2.4em, itemsep=0.2em]
\item 错误
\item 正确
\end{enumerate}
\end{question}
\begin{answer}
B（正确）
\end{answer}
\begin{explanation}
陀螺仪角速度对时间积分可以得到角度变化，因此可解析姿态角，但实际应用中会有漂移误差。
\end{explanation}

\begin{question}{102. PCF8591 模块的 AD 的数值范围是 0–255。}
\begin{enumerate}[label=\Alph*., leftmargin=2.4em, itemsep=0.2em]
\item 错误
\item 正确
\end{enumerate}
\end{question}
\begin{answer}
B（正确）
\end{answer}
\begin{explanation}
PCF8591 是 8 位 A/D，数字输出范围为 0-255。
\end{explanation}

\begin{question}{103. PCF8591 的功能包括多路模拟输入、内置跟踪保持，PCF8591 的最大转化速率由 I2C 总线的最大速率决定的。}
\begin{enumerate}[label=\Alph*., leftmargin=2.4em, itemsep=0.2em]
\item 错误
\item 正确
\end{enumerate}
\end{question}
\begin{answer}
B（正确）
\end{answer}
\begin{explanation}
PCF8591 通过 I2C 读写数据，转换与传输速率受 I2C 总线速率限制。
\end{explanation}

\begin{question}{104. 人脸检测指对于任意一幅给定的图像/视频帧，采用一定的策略对其进行搜索以确定 其中是否含有人脸，如果是则返回人脸的位置、大小和姿态。}
\begin{enumerate}[label=\Alph*., leftmargin=2.4em, itemsep=0.2em]
\item 错误
\item 正确
\end{enumerate}
\end{question}
\begin{answer}
B（正确）
\end{answer}
\begin{explanation}
人脸检测是在图像或视频帧中搜索人脸区域，并返回位置、大小等信息。
\end{explanation}

\begin{question}{105. MPU6050 陀螺仪传感器与树莓派通讯时使用的是 I2C 接口。}
\begin{enumerate}[label=\Alph*., leftmargin=2.4em, itemsep=0.2em]
\item 正确
\item 错误
\end{enumerate}
\end{question}
\begin{answer}
A（正确）
\end{answer}
\begin{explanation}
MPU6050 常通过 I2C 与树莓派通信，使用 SDA、SCL 两根主要信号线。
\end{explanation}

\begin{question}{106. LED 的导通压降为 0.3V。}
\begin{enumerate}[label=\Alph*., leftmargin=2.4em, itemsep=0.2em]
\item 正确
\item 错误
\end{enumerate}
\end{question}
\begin{answer}
B（错误）
\end{answer}
\begin{explanation}
普通 LED 的正向压降通常约 1.8V-3.3V，随颜色而变，0.3V 明显不正确。
\end{explanation}

\begin{question}{107. 双色 LED 中串联的 220 Ω 电阻的作用是限流，以免 LED 烧坏。}
\begin{enumerate}[label=\Alph*., leftmargin=2.4em, itemsep=0.2em]
\item 错误
\item 正确
\end{enumerate}
\end{question}
\begin{answer}
B（正确）
\end{answer}
\begin{explanation}
串联 220欧姆 电阻用于限制 LED 电流，防止电流过大烧坏 LED 或 GPIO。
\end{explanation}

\begin{question}{108. 电阻的选型中，除了电阻值和精度之外，还需要考虑电阻的功率。}
\begin{enumerate}[label=\Alph*., leftmargin=2.4em, itemsep=0.2em]
\item 正确
\item 错误
\end{enumerate}
\end{question}
\begin{answer}
A（正确）
\end{answer}
\begin{explanation}
电阻选型除阻值、精度外，还要考虑额定功率，否则可能过热损坏。
\end{explanation}

\begin{question}{109. 贴片钽电容正面深颜色标出的是负极，另一端为正极。}
\begin{enumerate}[label=\Alph*., leftmargin=2.4em, itemsep=0.2em]
\item 正确
\item 错误
\end{enumerate}
\end{question}
\begin{answer}
B（错误）
\end{answer}
\begin{explanation}
贴片钽电容通常用色带或标记表示正极，题干说深色标出负极不正确。
\end{explanation}

\begin{question}{110. 肖特基二极管是由金属-半导体组成的，具有超高速、低功耗等优点。}
\begin{enumerate}[label=\Alph*., leftmargin=2.4em, itemsep=0.2em]
\item 正确
\item 错误
\end{enumerate}
\end{question}
\begin{answer}
A（正确）
\end{answer}
\begin{explanation}
肖特基二极管由金属和半导体接触形成，正向压降低、开关速度快、功耗低。
\end{explanation}

\begin{question}{111. LED 的导通压降和普通二极管相同，约为 0.7V。}
\begin{enumerate}[label=\Alph*., leftmargin=2.4em, itemsep=0.2em]
\item 正确
\item 错误
\end{enumerate}
\end{question}
\begin{answer}
B（错误）
\end{answer}
\begin{explanation}
LED 正向压降通常高于普通硅二极管 0.7V，且不同颜色差异明显，因此说相同错误。
\end{explanation}

\begin{question}{112. 焊接中，提高焊接温度可以增大表面张力。}
\begin{enumerate}[label=\Alph*., leftmargin=2.4em, itemsep=0.2em]
\item 正确
\item 错误
\end{enumerate}
\end{question}
\begin{answer}
B（错误）
\end{answer}
\begin{explanation}
提高焊接温度通常会降低焊料表面张力、改善润湿，不是增大表面张力。
\end{explanation}

\begin{question}{113. 元器件管脚和焊盘的距离增大，可以使金属更好地扩散。}
\begin{enumerate}[label=\Alph*., leftmargin=2.4em, itemsep=0.2em]
\item 正确
\item 错误
\end{enumerate}
\end{question}
\begin{answer}
B（错误）
\end{answer}
\begin{explanation}
管脚与焊盘距离过大会影响润湿和焊点形成，不利于形成可靠结合层。
\end{explanation}

\begin{question}{114. 焊接界面的结合层最佳厚度为 1.2–3.5um。}
\begin{enumerate}[label=\Alph*., leftmargin=2.4em, itemsep=0.2em]
\item 正确
\item 错误
\end{enumerate}
\end{question}
\begin{answer}
A（正确）
\end{answer}
\begin{explanation}
焊接界面结合层过薄或过厚都不好，1.2-3.5um 属于常见的较佳范围。
\end{explanation}

\begin{question}{115. 内热式电烙铁，其发热元件在电烙铁头的内部。}
\begin{enumerate}[label=\Alph*., leftmargin=2.4em, itemsep=0.2em]
\item 正确
\item 错误
\end{enumerate}
\end{question}
\begin{answer}
A（正确）
\end{answer}
\begin{explanation}
内热式电烙铁的加热芯位于烙铁头内部或靠近内部，热效率较高。
\end{explanation}

\begin{question}{116. 手工焊接时，应该先撤离电烙铁再撤离焊锡。}
\begin{enumerate}[label=\Alph*., leftmargin=2.4em, itemsep=0.2em]
\item 正确
\item 错误
\end{enumerate}
\end{question}
\begin{answer}
B（错误）
\end{answer}
\begin{explanation}
手工焊接收尾应先撤离焊锡丝，再撤离电烙铁，使焊点继续受热润湿并自然凝固。
\end{explanation}

\begin{question}{117. 贴片元器件的焊接过程为：清洁 → 上锡 → 固定 → 焊接 → 清理。}
\begin{enumerate}[label=\Alph*., leftmargin=2.4em, itemsep=0.2em]
\item 正确
\item 错误
\end{enumerate}
\end{question}
\begin{answer}
A（正确）
\end{answer}
\begin{explanation}
贴片元器件手工焊接一般包括清洁、上锡、固定、焊接和清理等步骤。
\end{explanation}

\begin{question}{118. 贴片元器件的手工焊接过程中，小镊子的主要作用是夹取焊锡。}
\begin{enumerate}[label=\Alph*., leftmargin=2.4em, itemsep=0.2em]
\item 正确
\item 错误
\end{enumerate}
\end{question}
\begin{answer}
B（错误）
\end{answer}
\begin{explanation}
小镊子主要用来夹取和固定贴片元件，不是夹取焊锡。
\end{explanation}

\begin{question}{119. 润湿角越大，润湿力越大。}
\begin{enumerate}[label=\Alph*., leftmargin=2.4em, itemsep=0.2em]
\item 正确
\item 错误
\end{enumerate}
\end{question}
\begin{answer}
B（错误）
\end{answer}
\begin{explanation}
润湿角越小表示润湿越好，润湿力越大；润湿角越大反而润湿差。
\end{explanation}

\begin{question}{120. 贴片机的工作流程是：锡膏印刷 → PCB 传送至指定位置 → 贴装头取元器件 → 贴装。}
\begin{enumerate}[label=\Alph*., leftmargin=2.4em, itemsep=0.2em]
\item 正确
\item 错误
\end{enumerate}
\end{question}
\begin{answer}
A（正确）
\end{answer}
\begin{explanation}
贴片装配流程通常为锡膏印刷、PCB 传送定位、贴装头取料、贴装。
\end{explanation}

\begin{question}{121. MQTT 协议采用的是订阅、发布机制。}
\begin{enumerate}[label=\Alph*., leftmargin=2.4em, itemsep=0.2em]
\item 正确
\item 错误
\end{enumerate}
\end{question}
\begin{answer}
A（正确）
\end{answer}
\begin{explanation}
MQTT 采用发布/订阅机制，设备向主题发布消息或订阅主题接收消息。
\end{explanation}

\begin{question}{122. 在树莓派系统中，实时显示各进程资源使用情况的命令是 top。}
\begin{enumerate}[label=\Alph*., leftmargin=2.4em, itemsep=0.2em]
\item 正确
\item 错误
\end{enumerate}
\end{question}
\begin{answer}
A（正确）
\end{answer}
\begin{explanation}
top 命令可实时显示进程 CPU、内存等资源占用情况。 四、填空题
\end{explanation}

\section{四、填空题}

\begin{question}{123. 树莓派编程中，程序运行完成后，使用 \_\_ 按键退出程序。}
\end{question}
\begin{answer}
Ctrl+C
\end{answer}
\begin{explanation}
在终端运行程序时，Ctrl+C 用于发送中断信号，常用于退出正在运行的程序。
\end{explanation}

\begin{question}{124. 采样是将时间上的（a. 连续变化，b. 间隔变化）的模拟量转换成时间上（a. 连续变 化，b. 间隔变化）的模拟量。}
\end{question}
\begin{answer}
a b
\end{answer}
\begin{explanation}
采样是把时间上连续变化的模拟量变为时间上离散/间隔变化的模拟量。
\end{explanation}

\begin{question}{125. PCF8591 是一个单片集成、单独供电、低功耗、\_\_ 位 CMOS 数据获取器件。 PCF8591 其中具有 \_\_ 个模拟输入、\_\_ 个模拟输出和 \_\_ 个串行 I2C 总线接口。}
\end{question}
\begin{answer}
8 4 1 1
\end{answer}
\begin{explanation}
PCF8591 是 8 位器件，具有 4 个模拟输入、1 个模拟输出和 1 个串行 I2C 总线接口。
\end{explanation}

\begin{question}{126. 第一次使用树莓派时，需要在终端输入命令 \_\_\_\_ 以进入配置界面，进行一些简单 的配置。}
\end{question}
\begin{answer}
sudo raspi-config
\end{answer}
\begin{explanation}
sudo raspi-config 是树莓派系统常用配置命令，可进入配置界面修改接口、地区、密码等设置。
\end{explanation}

\begin{question}{127. 树莓派 4B 有 \_\_\_\_Pin 的 GPIO 接口。}
\end{question}
\begin{answer}
40
\end{answer}
\begin{explanation}
树莓派 4B 的 GPIO 排针为 40Pin。
\end{explanation}

\begin{question}{128. PCF8591 采用典型的 I2C 总线接口器件寻址方法，即总线地址由 \_\_、\_\_ 和 \_\_ 组成。}
\end{question}
\begin{answer}
器件地址 引脚地址 方向位
\end{answer}
\begin{explanation}
PCF8591 的 I2C 地址由器件地址、引脚地址和读写方向位共同组成。
\end{explanation}

\begin{question}{129. A/D 转换器是把输入的 \_\_ 转换成成比例的 \_\_，D/A 转换器是把输入 \_\_ 的 转换成成比例的 \_\_。}
\end{question}
\begin{answer}
模拟量 数字量 数字量 模拟量
\end{answer}
\begin{explanation}
A/D 把模拟量转换为数字量；D/A 把数字量转换为模拟量。
\end{explanation}

\begin{question}{130. 我们所用到的 PCF8591 的物理地址是 \_\_。}
\end{question}
\begin{answer}
0x48
\end{answer}
\begin{explanation}
常见 PCF8591 模块默认地址为 0x48，即 A0、A1、A2 地址脚为默认低电平时的地址。
\end{explanation}
